% GENERATED FILE. Edit canonical lessons, then run npm run generate:books.
% canonical-source-hash: fnv1a64:ef65447c5c059248
% canonical-lessons: GU-C450-vahal, GU-C450-sneh, GU-C450-mamta, GU-C450-mohvu, GU-C450-lad

\chapter{Words of Love}
\label{ch:gu-words-of-love}
\hlchaptermodality{450}

\begin{chapteropening}
\textbf{By the end of this chapter:} \emph{I can talk about love in a family.}

Complete the last lesson of chapter 450: I can talk about love in a family.
\end{chapteropening}

\section[\texorpdfstring{\gu{વહાલ} (vahāl)}{વહાલ (vahāl)}]{\gu{વહાલ} (vahāl) — fondness; tender affection}

\label{lesson:GU-C450-vahal}



\begin{quote}
Before the new one: say the Gujarati for to be won over, then the Gujarati for pleased.
\end{quote}

\subsection*{You'll want to know: \gu{વહાલ}}
\textbf{\gu{વહાલ}} — \emph{vahāl} — \textquotedblleft{}fondness; tender affection\textquotedblright{}.

\begin{quote}
\gu{દાદીને} \gu{નાના} \gu{બાળક} \gu{પર} \gu{ઘણું} \gu{વહાલ} \gu{છે}.
\end{quote}

\begin{quote}
\emph{dādīne nānā bāḷak par ghaṇũ vahāl che.}
\end{quote}

\begin{quote}
Grandmother is very fond of the little child.
\end{quote}

A neuter noun.

\begin{cousinweb}
\textbf{\gu{વહાલ}} (\emph{vahāl}) is inherited from Sanskrit \emph{vallabha}, \textquotedblleft{}dear, desired\textquotedblright{}; CDIAL lists the Gujarati word, \textquotedblleft{}affection\textquotedblright{} (CDIAL 11427).
\end{cousinweb}

\begin{grammarlens}[title={What you've built}]
One more word for describing what you lived and felt.
\end{grammarlens}

\subsection*{Your turn}
\begin{itemize}
\raggedright
  \item \emph{Say it:} \emph{vahāl}
  \item \emph{Say it:} \emph{vahāl}, once more
  \item \emph{Say it:} \emph{dādīne nānā bāḷak par ghaṇũ vahāl che.} — Grandmother is very fond of the little child.
  \item \emph{Recall:} say the Gujarati for to be won over, then the Gujarati for pleased, then say \emph{vahāl} again
\end{itemize}

\begin{tcolorbox}[breakable,colback=teal!4,colframe=teal!35!black,title={Before you move on}]
What does \gu{વહાલ} mean? (\textquotedblleft{}Fondness; tender affection\textquotedblright{}.) Say it once more.
\end{tcolorbox}
\section[\texorpdfstring{\gu{સ્નેહ} (sneh)}{સ્નેહ (sneh)}]{\gu{સ્નેહ} (sneh) — affection; loving warmth}

\label{lesson:GU-C450-sneh}



\begin{quote}
Before the new one: say the Gujarati for pleased, then the Gujarati for fondness.
\end{quote}

\subsection*{You'll want to know: \gu{સ્નેહ}}
\textbf{\gu{સ્નેહ}} — \emph{sneh} — \textquotedblleft{}affection; loving warmth\textquotedblright{}.

\begin{quote}
\gu{બંને} \gu{બહેનો} \gu{વચ્ચે} \gu{ઘણો} \gu{સ્નેહ} \gu{છે}.
\end{quote}

\begin{quote}
\emph{banne baheno vacce ghaṇo sneh che.}
\end{quote}

\begin{quote}
There is great affection between the two sisters.
\end{quote}

A masculine noun.

\begin{cousinweb}
\textbf{\gu{સ્નેહ}} (\emph{sneh}) is a tatsama, Sanskrit \emph{sneha}, which Monier-Williams gives first as \textquotedblleft{}oiliness\textquotedblright{} and then as tenderness and love; the older spoken form \emph{neh} comes from the same word (CDIAL 13802).
\end{cousinweb}

\begin{grammarlens}[title={What you've built}]
One more word for describing what you lived and felt.
\end{grammarlens}

\subsection*{Your turn}
\begin{itemize}
\raggedright
  \item \emph{Say it:} \emph{sneh}
  \item \emph{Say it:} \emph{sneh}, once more
  \item \emph{Say it:} \emph{banne baheno vacce ghaṇo sneh che.} — There is great affection between the two sisters.
  \item \emph{Recall:} say the Gujarati for pleased, then the Gujarati for fondness, then say \emph{sneh} again
\end{itemize}

\begin{tcolorbox}[breakable,colback=teal!4,colframe=teal!35!black,title={Before you move on}]
What does \gu{સ્નેહ} mean? (\textquotedblleft{}Affection; loving warmth\textquotedblright{}.) Say it once more.
\end{tcolorbox}
\section[\texorpdfstring{\gu{મમતા} (mamtā)}{મમતા (mamtā)}]{\gu{મમતા} (mamtā) — a mother's love}

\label{lesson:GU-C450-mamta}



\begin{quote}
Before the new one: say the Gujarati for fondness, then the Gujarati for affection.
\end{quote}

\subsection*{You'll want to know: \gu{મમતા}}
\textbf{\gu{મમતા}} — \emph{mamtā} — \textquotedblleft{}a mother's love\textquotedblright{}.

\begin{quote}
\gu{માતાની} \gu{મમતા} \gu{સૌથી} \gu{મોટી} \gu{છે}.
\end{quote}

\begin{quote}
\emph{mātānī mamtā sauthī moṭī che.}
\end{quote}

\begin{quote}
A mother's love is the greatest of all.
\end{quote}

A feminine noun.

\begin{cousinweb}
\textbf{\gu{મમતા}} (\emph{mamtā}) is a tatsama: Monier-Williams gives Sanskrit \emph{mamatā} as the state of feeling \textquotedblleft{}mine\textquotedblright{}, from \emph{mama}, \textquotedblleft{}my\textquotedblright{}; Gujarati uses it for a mother's tender love.
\end{cousinweb}

\begin{grammarlens}[title={What you've built}]
One more word for describing what you lived and felt.
\end{grammarlens}

\subsection*{Your turn}
\begin{itemize}
\raggedright
  \item \emph{Say it:} \emph{mamtā}
  \item \emph{Say it:} \emph{mamtā}, once more
  \item \emph{Say it:} \emph{mātānī mamtā sauthī moṭī che.} — A mother's love is the greatest of all.
  \item \emph{Recall:} say the Gujarati for fondness, then the Gujarati for affection, then say \emph{mamtā} again
\end{itemize}

\begin{tcolorbox}[breakable,colback=teal!4,colframe=teal!35!black,title={Before you move on}]
What does \gu{મમતા} mean? (\textquotedblleft{}A mother's love\textquotedblright{}.) Say it once more.
\end{tcolorbox}
\section[\texorpdfstring{\gu{મોહવું} (mohvũ)}{મોહવું (mohvũ)}]{\gu{મોહવું} (mohvũ) — to charm; to captivate}

\label{lesson:GU-C450-mohvu}



\begin{quote}
Before the new one: say the Gujarati for affection, then the Gujarati for a mother's love.
\end{quote}

\subsection*{You'll want to know: \gu{મોહવું}}
\textbf{\gu{મોહવું}} — \emph{mohvũ} — \textquotedblleft{}to charm; to captivate\textquotedblright{}.

\begin{quote}
\gu{નાના} \gu{બાળકનું} \gu{સ્મિત} \gu{બધાને} \gu{મોહે} \gu{છે}.
\end{quote}

\begin{quote}
\emph{nānā bāḷaknũ smit badhāne mohe che.}
\end{quote}

\begin{quote}
A little child's smile charms everyone.
\end{quote}

Past, with \gu{તેણે} (he or she …): \gu{તેણે} \gu{બધાનું} \gu{હૃદય} \gu{મોહ્યું}, \emph{teṇe badhānũ hṛday mohyũ}.

\begin{cousinweb}
\textbf{\gu{મોહવું}} (\emph{mohvũ}) is inherited from Sanskrit \emph{mōhayati}, \textquotedblleft{}bewilders\textquotedblright{}; CDIAL lists the Gujarati verb, \textquotedblleft{}to allure, charm\textquotedblright{} (CDIAL 10362).
\end{cousinweb}

\begin{grammarlens}[title={What you've built}]
One more word for describing what you lived and felt.
\end{grammarlens}

\subsection*{Your turn}
\begin{itemize}
\raggedright
  \item \emph{Say it:} \emph{mohvũ}
  \item \emph{Say it:} \emph{mohvũ}, once more
  \item \emph{Say it:} \emph{nānā bāḷaknũ smit badhāne mohe che.} — A little child's smile charms everyone.
  \item \emph{Recall:} say the Gujarati for affection, then the Gujarati for a mother's love, then say \emph{mohvũ} again
\end{itemize}

\begin{tcolorbox}[breakable,colback=teal!4,colframe=teal!35!black,title={Before you move on}]
What does \gu{મોહવું} mean? (\textquotedblleft{}To charm; to captivate\textquotedblright{}.) Say it once more.
\end{tcolorbox}
\section[\texorpdfstring{\gu{લાડ} (lāḍ)}{લાડ (lāḍ)}]{\gu{લાડ} (lāḍ) — pampering; fond indulgence}

\label{lesson:GU-C450-lad}



\begin{quote}
Before the new one: say the Gujarati for a mother's love, then the Gujarati for to captivate.
\end{quote}

\subsection*{You'll want to know: \gu{લાડ}}
\textbf{\gu{લાડ}} — \emph{lāḍ} — \textquotedblleft{}pampering; fond indulgence\textquotedblright{}.

\begin{quote}
\gu{દાદાના} \gu{લાડથી} \gu{બાળક} \gu{જિદ્દી} \gu{થઈ} \gu{ગયું}.
\end{quote}

\begin{quote}
\emph{dādānā lāḍthī bāḷak jiddī thaī gayũ.}
\end{quote}

\begin{quote}
Grandfather's pampering has made the child stubborn.
\end{quote}

A neuter noun.

\begin{cousinweb}
\textbf{\gu{લાડ}} (\emph{lāḍ}) goes back to a reconstructed \emph{*lāḍya-}, \textquotedblleft{}fondling\textquotedblright{}; CDIAL lists the Gujarati word, \textquotedblleft{}caress\textquotedblright{} (CDIAL 11013).
\end{cousinweb}

\begin{grammarlens}[title={What you've built}]
One more word for describing what you lived and felt.
\end{grammarlens}

\subsection*{Your turn}
\begin{itemize}
\raggedright
  \item \emph{Say it:} \emph{lāḍ}
  \item \emph{Say it:} \emph{lāḍ}, once more
  \item \emph{Say it:} \emph{dādānā lāḍthī bāḷak jiddī thaī gayũ.} — Grandfather's pampering has made the child stubborn.
  \item \emph{Recall:} say the Gujarati for a mother's love, then the Gujarati for to captivate, then say \emph{lāḍ} again
\end{itemize}

\begin{tcolorbox}[breakable,colback=teal!4,colframe=teal!35!black,title={Before you move on}]
What does \gu{લાડ} mean? (\textquotedblleft{}Pampering; fond indulgence\textquotedblright{}.) Say it once more.
\end{tcolorbox}
